\documentclass[a4paper]{article}
\usepackage[pdftex]{graphicx}
\usepackage[utf8]{inputenc}
\usepackage{enumerate}
\usepackage{siunitx}
\sisetup{locale=DE} 
\usepackage{href-ul}
\hypersetup{
	colorlinks=true,
	linkcolor=blue,
	urlcolor=blue}
\usepackage{icomma}
\usepackage{amssymb}
\usepackage{geometry}
\geometry{a4paper, top=15mm, left=15mm, right=15mm, bottom=15mm,
	headsep=10mm, footskip=12mm}

\begin{document}

\begin{enumerate}[1.]

\item {\bf Schwingungen; Fadenpendel}\\
Bei einem Pendel hat man den Zusammenhang zwischen Pendellänge und Frequenz gemessen:

\begin{tabular}{cccccc}
$l$ in m & 0,20 & 0,40 & 0,60 & 0,80 & 1,0 \\
$f$ in Hz & 1,1 & 0,79 & 0,64 & 0,56 & 0,50  
\end{tabular}

\begin{enumerate}[a)]
\item Zeichne ein Diagramm, aus dem der Zusammenhang zwischen der Länge $l$ und der \underline{Schwingungsdauer} $T$ abgelesen werden kann.\\
(x-Achse= $ ~ l $-Achse;  $\SI{10}{\centi\meter}~\hat{=}~\SI{1}{\meter}$ ;
 y-Achse=$ ~T$-Achse: $\SI{5}{\centi\meter}~\hat{=}~\SI{1}{\second}$)
\item Wie lang ist ein Pendel, das eine Schwingungsdauer von $\SI{1,5}{\second}$ hat ? Benutze dein Diagramm.
\item Von welchen Eigenschaften des Pendels ist die Schwingungsdauer abhängig?
\item Ein Pendel besitzt eine Frequenz von $f=\SI{0,5}{\hertz}$. Wie viele Perioden vergehen in $\SI{24}{\second}$? 
\end{enumerate}


{\bf \item Die schwingende Masse eines Federpendels beträgt $\SI{440}{\gram}$.} Die Schwingung erfolgt mit einer Maximalgeschwindigkeit des Pendelkörpers von $\SI{0,30}{\meter\over\second}$. Beim Versuch werden für 10 Schwingungen $\SI{12}{\second}$ benötigt. Die Federhärte beträgt $\SI{12}{\newton\over\meter}$. Zum Zeitpunkt $t = \SI{0}{\second}$ schwingt das Federpendel in Abwärts\-richtung durch die Ruhelage 

\begin{enumerate}[a)]
\item  Wann bezeichnet man eine Schwingung als harmonische Schwingung? Gib 2 Charakterisierungen dafür an.
\item Die Schwingungsdauer des Systems lässt sich auf zwei Arten aus den Angaben berechnen. Vergleiche die Ergebnisse und gib die Frequenz an.
\item Stelle die Funktion $v(t)$ für die Geschwindigkeit des Pendelkörpers auf und berechne sie zur Zeit $ t=\SI{7,9}{\second}$. \\
Skizziere den Verlauf des Graphen von $v(t)$ im Intervall $[0;\SI{1,2}{\second}]$. Markiere durch farbige Pfeile die Zeitpunkte der größten Beschleunigungsbeträge.
\item Erläutere den Ansatz, durch den wir im Unterricht einen Zusammenhang zwi\-schen der Amplitude und der Maximalgeschwindigkeit hergestellt haben und  berechne damit die Amplitude der Schwingung.\\
$[\textnormal{Falsches Ersatzergebnis}: \hat{x}=\SI{5,0}{\centi\meter}]$
\item Gib Größe und Richtung der rücktreibenden Kraft im oberen Umkehrpunkt an.

\item Durch einen Austausch der Feder bei gleicher Masse erhöht sich die Schwing\-ungsdauer auf $T=\SI{1,5}{\second}$. \\
Welche Federhärte besitzt die neue Feder?
\end{enumerate} 

{\bf \item Gib jeweils eine Veränderung an einem Federpendel an, durch die...}

\begin{enumerate}[a)]
\item ...sich die Schwingungsdauer verdoppelt
\item ...die rücktreibende Kraft bei gleicher Auslenkung nur noch halb so groß ist
\end{enumerate}  

\end{enumerate} 

\fbox{
	\begin{minipage}{0.5\textwidth}
		Zur Lösung bitte \href{https://www.okuyakl.de/phys/p10penL407/ll407.pdf}{hier klicken} oder den QR-Code scannen.\\
	Weitere Arbeitsblätter gibt es unter 
	
	\href{https://www.okuyakl.de}{www.okuyakl.de}
	\end{minipage}
	\hfill
	\begin{minipage}{0.4\textwidth}
		\includegraphics[width=1.5 cm]{../../viecher/zwe03}
		\includegraphics[width=3 cm]{qr407}
		\includegraphics[width=2 cm]{../../viecher/afanticon1}
		
	\end{minipage}}

\end{document}